\begin{enumerate}[label=\textbf{\ansnum.\arabic*}]
  \item $49\cdot10=490$ bytes/s, about $3.9$\,kbit/s. At $50$ frames per second: $19.6$\,kbit/s per robot, so $5\,000/19.6\approx255$ robots; the debug stream is tiny. (Images would be a different story.)
  \item The two computers' clocks are not synchronised, and \verb!time.monotonic()! on two computers has completely different zero points. A command could look $1000$\,s old, or valid for ever. The robot stamps commands with its own clock when they arrive.
  \item
\begin{verbatim}
import queue

def __init__(self, slot, port=5555):
    ...
    self.out = queue.Queue(maxsize=1)
    threading.Thread(target=self._send, daemon=True).start()

def publish(self, frame):
    try:
        self.out.get_nowait()          # drop the old frame if unsent
    except queue.Empty:
        pass
    self.out.put_nowait(frame)

def _send(self):
    while True:
        frame = self.out.get()
        conn = self.conn
        if conn is not None:
            try:
                send_msg(conn, frame)
            except OSError:
                self.conn = None
\end{verbatim}
(With one producer, the get-then-put pair is safe.)
  \item Yes in both directions. An old robotcontrol skips the unknown field number $8$ when parsing. An old robot does not send it, and the new robotcontrol sees an empty list. That is the reason for field numbers in protobuf.
  \item Listen only on the team's own network interface or address (not \texttt{0.0.0.0}), and require a shared team token in the first message of a connection before commands are accepted; also accept commands only while the robot is not in a game (the GameController state says so).
\end{enumerate}
